\documentclass{ximera}
\title{Inverses, Summary}

\newcommand{\pskip}{\vskip 0.1 in}

\begin{document}
\begin{abstract}
Summary of inverse functions.
\end{abstract}
\maketitle

\section{The Basics}

\begin{question}  \label{Qdf3r3fd44xxz}
Let 
\[
     y =f(x) = \frac{6}{7}x + 4.
\]

\begin{enumerate}
\item Describe the action of $f$ on its input.

\item Describe the action of $f^{-1}$ on its input.

\item Use your description in part (b) to find an expression for the function $x=f^{-1}(y)$.

\item Verify algebraically that $f^{-1}$ undoes $f$.

\item Verify algebraically that $f$ undoes $f^{-1}$.
\end{enumerate}

\begin{explanation}
\begin{enumerate}
\item The function $f$ first multiplies its input by $6/7$ and then adds $4$ to the result.

\item The function $f^{-1}$ first subtracts $4$ from its input and then multiplies the result by $7/6$.

\item 
\[
   x = f^{-1}(y) = \frac{7}{6} \left( y - 4\right).
\]

\item To check that $f^{-1}$ undoes $f$ we must show that
\[
     f^{-1}(f(x)) = x
\]
for all $x$ in the domain of $f$.

Here's the algebra:

\begin{align*}
     f^{-1}(f(x)) &= f^{-1}\left( \frac{6}{7}x + 4  \right) \\ 
                       &= \frac{7}{6} \left( \left( \frac{6}{7}x + 4 \right) - 4 \right)  \\
                       &= \frac{7}{6} \left(   \frac{6}{7}x\right) \\
                       &= x.
\end{align*}

\item To check that $f$ undoes $f^{-1}$ we must show that
\[
     f(f^{-1}(y)) = y
\]
for all $y$ in the domain of $f$.

Here's the algebra:

\begin{align*}
     f(f^{-1}(y)) &= f\left(   \frac{7}{6} \left( y - 4\right)\right) \\ 
                       &= \frac{6}{7} \left( \frac{7}{6} \left( y - 4\right) \right) +4  \\
                       &= (y-4) + 4 \\
                       &= y.
\end{align*}


\end{enumerate}
\end{explanation}
\end{question}

\section{One-to-One Functions}

The basic idea is this. The inverse of a function is always a relation (as set of ordered pairs), but it need not be a function.

For example, the function 
\[
   y=f(x) = x^2 , x\in \mathbb{R},
\]
is really the set of ordered pairs
\[
  {\cal R} = \{  (x,x^2) \, | \, x \in \mathbb{R} \}. 
\]
This relation is a function because each input has \emph{exactly} one output. The graph of the relation passes the horizontal line test - ie. each horizontal line intersect the graph of the function $y=f(x)=x^2$ in at most one point. The test is a way to visualize the fact that each input (ie. $x$-coordinate) has exactly one output (ie. $y$-coordinate).

We get the inverse of any relation by swapping the $x$ and $y$-coordinates. The inverse of the above function, for example, is the relation
\[
  {\cal R}^{-1}  = \{  (x^2, x)  \, | \, x \in \mathbb{R}\}.
\] 

This relation is \emph{not} a function since there is at least one input that has more than one output. For example, the points $(9,3)$ and $(9,-3)$ are both in the relation ${\cal R}^{-1}$, meaning that in this relation the input $9$ has two outputs, $3$ and $-3$.

Geometrically the graph of ${\cal R}^{-1}$ ($x=y^2$) fails the vertical line test - ie. the vertical line $x=9$ intersects the graph in more than one point. This is because the graph of the original relation ($y=x^2$) fails the horizontal line test. That is, the horizontal line $y=9$ intersects the graph of $y=x^2$ in more than one point. And this is because the function
\[
 y = f(x) = x^2 \, , \, x\in \mathbb{R}
\]
is not \emph{one-to-one}. 

What is means for a function to be \emph{one-to-one} is that \emph{any two different inputs always have different outputs}. And a function is not one-to-one when there is at least one pair of different inputs with the same outputs.

The function $y=f(x) = x^2$, for example, is not one-to-one because $f(3) = f(-3)$.

Here's the simple way to look at all this. I pick a number, square it, and get $25$. What was my number? Well, there are two possibilities, my number could have been $5$ or $-5$. You can't say for sure. So the squaring function is not one-to-one and thefore it's inverse is not a function. 

On the other hand, the function
\[
   y = g(x) = x^2 \, , \, x \leq 0,
\]
 is one-to-one. If I now pick a number less than or equal to zero, square it, and get $25$, for example, there is only one possibility. My number must have been $-5$. And the same can be said if I pick any non-positive number. If I tell you the square of my number, you know can find my original number for sure.
 
\begin{example} \label{Ev77337}
Find the inverse of the function 
\[
   y = g(x) = x^2 \, , \, x \leq 0.
\]

\begin{explanation}
Since
\[
   y = x^2 ,
\]
\[
    x = \pm \sqrt{y}.
\]
But since $x\leq 0$, 
\[
   x = g^{-1}(y) = -\sqrt{y} \, , \, y \geq 0.
\]
\end{explanation}
\end{example}

\begin{example}     \label{Ex99:fefQuadratics}
Let 
\[
   f(x) = \sqrt{25-x^2} \, , -5 \leq x \leq 5 .
\]

\begin{enumerate}
\item Use your knowledge of circles to sketch a graph of the function $y=f(x)$ by hand.

\item Explain what it means for a funtion to be one-to-one.

\item Is the function $f$ one-to-one? Explain your reasoning.

\item Explain what it means for a relation to be a function.

\item Is the inverse of $f$ a function? Explain your reasoning.

\end{enumerate}

\begin{explanation}
\begin{enumerate}
\item Since 
\[
    y = \sqrt{25-x^2},
\]
we know that $y \geq 0$. Then squaring both sides and doing some algebra (you should fill in the details) gives
\[
    x^2 + y^2 =25 \text{ and } y\geq 0.
\]
The curve 
\[
   x^2 + y^2 = 25
\]
is a circle of radius 5 centered at the origin. But since $y\geq 0$, the graph (not shown) of the function $y=f(x)$ is a semicircle in the first and second quadrants with endpoints $(5,0)$ and $(-5,0)$.

\item A function is one-to-one if any two distinct inputs have distinct outputs.

\item The function $f$ is {\bf not} one-to-one. We can see this because the two inputs $x=-5$ and $x=5$ give the same output
\[
     f(5) = \sqrt{25-5^2} = 0 
\]
and
\[
     f(-5) = \sqrt{25-(-5)^2} = 0 = f(5). 
\]

We can also see that $f$ is not one-to-one from the graph. There are many horizontal lines that intersect the graph in two points. This means that there are many pairs of inputs that give the same output.

\item A function is a relation when each input has exactly one output.

\item Because $f$ is not one-to-one its inverse is {\bf not} a function. For example, the input $y=0$ to the inverse relation would have two outputs, $x=-5$ or $x=5$, making the inverse relation not a function.

\end{enumerate}

\end{explanation}
\end{example}


\begin{example} \label{Ex98:5Quadratics}
Let 
\[
    y=   f(x) = \sqrt{25-x^2} \, , 0\leq x \leq 5 .
\]
\begin{enumerate}
\item Use your knowledge of circles to graph the function $y=f(x)$ \emph{by hand}. Use the graph to thoroughly explain why $f$ is one-to-one.

\item Find an expression for the inverse function
\[
    x = f^{-1}(y)
\]
that expresses the input of $f$ in terms of its output. State the domain of $f^{-1}$ using set notation.

\item Show algebraically that $f^{-1}$ undoes $f$.

\item Show algebraically that $f$ undoes $f^{-1}$.

\item Graph the function $y=f^{-1}(x)$. 

\item Explain how you could have realized that $f^{-1}(x) = f(x)$ by looking at the graph of $y=f(x)$.
\end{enumerate}


\begin{explanation}
\begin{enumerate}
\item The graph of $y=f(x)$ is a quarter-circle in the first quadrant centered at the origin with endpoints $(5,0)$ and $(0,5)$. Since each horizontal line intersects the graph in at most one point, any two different inputs to $f$ have different ouputs. This means that $f$ is one-to-one.

\item We know from the graph in part (a) that the range of $f$ is the set
\[
   \{ y \, | \,  0\leq y \leq 5 \} .
\]
This is the domain of the function $x=f^{-1}(y)$.

Now to find an expression for $f^{-1}(y)$, we solve the equation
\[
   y=   \sqrt{25-x^2} \, , 0\leq x \leq 5 
\]
for $x$. Squaring both sides gives
\[
   y^2 = 25 - x^2 .
\]
So 
\[
     x^2 = 25-y^2
\]
and 
\[
   x  = \pm \sqrt{25-y^2} .
\]

To decide which sign to choose, we go back to the domain of $f$, where we are given that 
\[
     0 \leq x \leq 5 .
\]
Since $x\geq 0$, this means we must choose the positive sign and so
\[
    x = f^{-1}(y) = \sqrt{25-y^2}.
\]
We should also include the domain $0\leq y \leq 5$ (as stated above). So the inverse of the function
\[
    y=   f(x) = \sqrt{25-x^2} \, , 0\leq x \leq 5 
\]
is the function
\[
    x = f^{-1}(y) = \sqrt{25-y^2} \, , 0\leq y \leq 5 .
\]

\item To show $f^{1-}$ undoes $f$, we simplify the composition $f^{-1}(f(x))$ as follows:
\begin{align*}
           f^{-1}(f(x)) &= f^{-1}(\sqrt{25-x^2})   \\
                             &= \sqrt{25 - \left( \sqrt{25-x^2}\right)^2}  \\
                             &= \sqrt{25 - (25-x^2)}  \\
                             &= \sqrt{x^2} \\
                             &= |x|  \\
                             &= x \text{ since } 0\leq x \leq 5.
\end{align*}

\item To show $f$ undoes $f^{-1}$ simplify the composition $f(f^{-1}(x))$. I'll leave this for you.

\begin{freeResponse}
\end{freeResponse}

\item Since the variables $x$ and $y$ have no practical meaning, we can switch them in our expression for $x = f^{-1}(y)$. So the inverse of the function
\[
    y=   f(x) = \sqrt{25-x^2} \, , 0\leq x \leq 5 
\]
is the function
\[
    y = f^{-1}(x) = \sqrt{25-x^2} \, , 0\leq x \leq 5 .
\]
This is identical to the function $f$. That is, $f$ is its own inverse and the graphs of $f$ and $f^{-1}$ are identical.

\item Since the graph of $y=f(x)$ is symmetric about the line $y=x$, we could have realized from the start that $f$ is its own inverse.

\end{enumerate}
\end{explanation}
\end{example}



\section{Applications}

\begin{example} \label{QlvemME8381O}
\item The function 
\[
  g = f(v)  \, , \, 25\leq v \leq 50 ,
\] 
expresses the gas mileage of a car (in miles/gal) in terms of its speed (in miles/hr).
   

\begin{enumerate}
   \item Explain what it would mean for the function $f$ to be one-to-one in this scenario.

  \item Is it likely that $f$ is one-to-one? Explain your reasoning.    
\end{enumerate}

\begin{explanation}
\begin{enumerate}
\item For the function $f$ to be one-to-one means the car would get different gas mileages for any pair of different speeds between $25$ and $50$ miles/hour.

\item You could say that $f$ is likely to be one-to-one because you think the function $f$ is increasing, i.e. the greater the speed, the greater the gas mileage (at least for speeds between $25$ miles/hour and $50$ miles/hour).

You could alsos say that $f$ is unlikely to be one-to-one because you think it likely the car gets the same gas mileage at two different speeds between $25$ miles/hour and $50$ miles/hour.

Either answer is correct. The important thing is your logic.
\end{enumerate}
\end{explanation}
\end{example}


\begin{example} \label{Emv3et3}
The function
\[
g = f(v) = 40 -\frac{1}{25}\left(v-50 \right)^2 \, , \, 25\leq v \leq 50, 
\]
expresses the gas mileage of a car (in miles/gal) in terms of its speed (in mile

\begin{enumerate}
\item Find an expression for the function $v=f^{-1}(g)$. Include the domain.
     \item Show algebraicall that $f^{-1}$ undoes $f$.
    \item Simplify the function
    \[
             G =h(g)= f(f^{-1}(g)-10).
    \]
    Include the domain and interpret its meaning.
    \end{enumerate}

\begin{explanation}
\begin{enumerate}
\item Since 
\[
  g =  40 -\frac{1}{25}\left(v-50 \right)^2 ,
\]
\[
  (v-50)^2 = 900 - 25g ,
\]
and
\[
 v  = 50 \pm \sqrt{900-25g} .
\]
To choose between $\pm$, we go back to the domain of $f$, where $25 \leq v \leq 50$. Since $\sqrt{900-25g}\geq 0$, we need to choose the minus sign to get speeds at most $50$ miles/hour. So
\[
 v = f^{-1}(g) = 50 - \sqrt{900-25g}.
\]

\item To show $f^{-1}$ undoes $f$ we need to show
\[
   f^{-1}(f(v)) = v \, , \, 25\leq v \leq 50.
\]

Here's the computation.

\begin{align*}
   f^{-1}(f(v)) &= f^{-1}\left(  40 -\frac{1}{25}\left(v-50 \right)^2 \right) \\ 
                     &= 50 - \sqrt{900 - 25 \left( 40 -\frac{1}{25}\left(v-50 \right)^2 \right)} \\ 
                     &= 50 - \sqrt{(v-50)^2} \\
                     &= 50 - \left| v-50 \right| .
\end{align*}

Now we stop and think.

To simplify $| v-50 |$, we need to know if $v-50$ is at least zero or at most zero. To decide, we go back to the domain of $f$, where we were told $25\leq v\leq 50$. This domain implies
\[
    25 - 50 \leq  v - 50 \leq 50 - 50 ,
\]
or equivalently that
\[
   -25 \leq v - 50 \leq 0.
\]
So $v - 50\leq 0$ and
\begin{align*}
  | v - 50 | &= - (v-50) \\
                & = 50 - v.
\end{align*}

Now we can continue simplifying the composition.
\begin{align*}
  f^{-1}(f(v)) &= 50 - \left| v-50 \right| \\
                    &= 50 - (50 -v) \\
                    &=v. 
\end{align*}


\end{enumerate}
\end{explanation}

\end{example}



\end{document}